\documentclass[letterpaper,twocolumn,10pt]{article}
\usepackage{hyperref}
\usepackage{usenix2019_v3_xetex}
\usepackage[T1]{fontenc}
\usepackage{amsmath,amssymb,graphicx}
\hypersetup{hidelinks,pdftitle={Architecture figure: typesetting proof},pdfauthor={}}
\makeatletter
\newcommand{\captiontypefigure}{\def\@captype{figure}}
\makeatother
\begin{document}
\twocolumn[{
\noindent\textbf{Architecture figure: author-review proof}\par
\vspace{8pt}
\centering
\includegraphics[width=\textwidth]{figures/system-architecture-technical.pdf}
\vspace{3pt}
\captiontypefigure
\caption{Shared state and per-object layouts.
(a) Two illustrative views change different ranges of one file.
$\Delta_i$ resides in the mutable metadata service; private payload $U_i$ and
snapshot $S$ reside on object storage. The index--record--frame path represents
a name lookup. Upper reads verify whole objects before slicing; visibility and
Hole rules are omitted here.
(b) Fixed policies partition the same 1\,MiB object on a common byte scale.
The enlarged 2+2 inset has $s=B/2=128\,\mathrm{KiB}$ and $q=64\,\mathrm{KiB}$;
the 1\,MiB policy instead has $s=q=512\,\mathrm{KiB}$.
These are supported configurations, not automatic-policy outcomes.
RustFS persists the admitted layout for reads and repair. Smaller-$q$ writes
require both default-off gates; dashed hint wiring still needs Entity/producer
integration. Packed frames and $q$ are distinct.}

\vspace{14pt}
}]
\section{Shared immutable state}
Let $S$ be a pinned immutable snapshot and $\Delta_i$ the private upper of
workspace $i$. Its logical view is $W_i=\operatorname{Resolve}(S,\Delta_i)$.
Private names, attributes and file mappings remain in the mutable metadata
service. The file-range example highlights separate private changes while
unchanged ranges retain explicit references to the same lower snapshot.

\section{Per-object storage layout}
The coding block $B$ is selected at the finalized-object boundary. The storage
backend admits and persists the effective geometry. Read and repair paths use
that stored geometry. The two equal-width strips compare supported fixed
layouts for the same illustrative object; they do not describe an adaptive
optimizer or a measured performance result.
\end{document}
